\section{Related Work}
\label{sec:relatedWork}

\noindent
{\bf Retrieval-Augmented Generation (RAG).}
Standard RAG~\cite{lewis2020rag} integrates external knowledge into the generation process to improve factual grounding. While subsequent hybrid variants~\cite{izacard2021fid} and persistent memory stores~\cite{park2023generative} improve scalability, they remain inherently \emph{retrospective}. These systems retrieve based on surface similarity to the current query rather than goal-oriented anticipation, often failing to surface memories that are only indirectly relevant to future needs.

\noindent
{\bf Graph-Based Retrieval and GraphRAG.}
To move beyond isolated text chunks, \textit{GraphRAG}~\cite{edge2024graphrag} leverages knowledge graphs to support multi-hop inference. However, these methods are constrained by a \emph{preordained structure}: retrieval is limited to traversing static edges established at construction time. Consequently, they cannot infer counterfactual relations or anticipate connections dynamically based on specific query -- a hallmark of human prospection.

\noindent
{\bf Agentic Memory and Reasoners.}
Recent frameworks such as \textit{ToT}~\cite{yao2023tree} and \textbf{ReAct}~\cite{yao2023react} use explicit reasoning traces to enhance planning and action. Agentic memory systems such as Mem0$^{\texttt{g}}$~\cite{chhikara2025mem0} maintain persistent user memories and entity relations, but retrieval still relies on vector search and extracted graph structure. In contrast, \textit{PGR} extends the principles of branching reasoning to the \emph{retrieval domain}: simulated ToT or chain steps become retrieval probes, and retrieved facts refine subsequent prospection. This transforms retrieval into an active, iterative policy that is orthogonal to—and can be layered upon—any underlying storage architecture.

\noindent
{\bf Neuroscience of Prospection.}
Cognitive neuroscience observes that human memory is constructive, not archival; hippocampal-prefrontal circuits support both remembering the past and imagining the future~\cite{schacter2017episodic, addis2007constructive}. This \emph{simulation} allows fragments of experience to be recombined for anticipatory planning~\cite{hassabis2007scenes}. PGR operationalizes this insight, aligning memory retrieval with active simulation mechanisms observed in human cognition.